% !TEX root = ../CV_Tramarin.tex

%%%%%%%%%%%%%%%%%%%%%%%%%%%%
%%%%%% ACADEMIC ACTIVITIES
%%%%%%%%%%%%%%%%%%%%%%%%%%%%

\sectioniten{Responsabilità e ruoli in ambito Accademico}{Responsibilities in Academic Activities}

\langit{
	\cventry{2026 --}{Commissione Ammissione Studenti Internazionali}{\newline{}Membro della Commissione per l'Ammissione degli studenti internazionali}{Corso di Laurea Magistrale in Mechatronics Engineering}{Dipartimento di Tecnica e Gestione dei Sistemi, Università di Padova.}{Con le seguenti competenze: \newline - Analizzare i profili dei candidati internazionali;\newline - Coadiuvare il Coordinatore del Corso di Laurea Magistrale nella definizione dei criteri di valutazione;\newline -  Coadiuvare il Coordinatore del Corso di Laurea Magistrale nella valutazione e selezione finale dei candidati.
	}	
	\cventry{2024 - 2025}{Commissione Assicurazione Qualità di Corso di Dottorato}{\newline{}Membro della Commissione per l'Assicurazione della Qualità del Corso di Dottorato in ICT}{}{Dipartimento di Ingegneria ``E. Ferrari'', Università di Modena e Reggio Emilia, Italia.}{Con le seguenti competenze: \newline - Analizzare i risultati delle indagini sottoposte ai dottorandi;\newline - Coadiuvare il Coordinatore del Corso di Dottorato nella verifica e stesura della Scheda Annuale;\newline -  Coadiuvare il Coordinatore del Corso di Dottorato nella stesura della Relazione di Monitoraggio del Corso di Dottorato.
	}	

	\cventry{2020 - 2025}{Commissione Assicurazione Qualità di Dipartimento}{\newline{}Membro della Commissione per l'Assicurazione della Qualità}{}{Dipartimento di Ingegneria ``E. Ferrari''}{Università di Modena e Reggio Emilia, Italia.}

	\cventry{2020 - 2025}{Commissione Assicurazione Qualità di CdS}{\newline{}Responsabile per l'Assicurazione della Qualità per i Corsi di Laurea Magistrale e Triennale in Ingegneria Elettronica}{}{Dipartimento di Ingegneria ``E. Ferrari'', Università di Modena e Reggio Emilia, Italia.}{Con le seguenti competenze: \newline - Analizzare i risultati delle indagini di soddisfazione degli studenti e dei laureati;\newline - Proporre azioni di miglioramento per il Corso di Studio in base ai risultati delle indagini di soddisfazione degli studenti e dei laureati;\newline - Analizzare i risultati e le osservazioni emerse nella Relazione della Commissione Paritetica Docenti Studenti;\newline - Proporre azioni di miglioramento per il Corso di Studio in base ai risultati osservazioni emerse nella Relazione della Commissione Paritetica Docenti Studenti;\newline - Coadiuvare il Presidente di CdS nella stesura della Scheda di Monitoraggio Annuale e del Rapporto di Riesame Ciclico;\newline - Coadiuvare il Presidente di CdS nella stesura della Relazione di Monitoraggio Annuale dell'Assicurazione Qualità del CdS.
	}

	\cventry{2020 - 2025}{Commissione Tirocini}{\newline{}Membro della Commissione per la Valutazione dei Tirocini}{}{Corso di Laurea Magistrale in Electric Vehicle Engineering, Università di Bologna, Italia.}{Con le seguenti competenze: \newline - Approvare i programmi di tirocinio;\newline - Valutare il tirocinio svolto per permettere l’acquisizione dei crediti formativi stabiliti dal regolamento del corso di studio.
	}

	\cventry{2020 - 2025}{Commissione Didattica}{\newline{}Membro della Commissione Didattica}{}{Corso di Laurea Magistrale in Electric Vehicle Engineering, Università di Bologna, Italia.}{Con le seguenti competenze: \newline - Deliberare tempestivamente - e comunque in tempo utile per garantire il rispetto dei termini procedimentali previsti dalla regolamentazione in materia - rispetto all’istanza formulata dallo studente in merito al riconoscimento dei crediti formativi acquisiti in altri corsi di studio dell’università o di altri atenei;\newline - Deliberare tempestivamente in merito all’analisi e validazione dei piani di studio personalizzati o individuali;\newline - Deliberare tempestivamente in merito al riconoscimento dei crediti da attività lavorativa, attività extra-universitaria o altra attività in sostituzione del tirocinio, alle competenze acquisite fuori dall’Università nei casi consentiti dal Regolamento didattico del corso di studio, nonché in merito ad ogni altro riconoscimento di crediti acquisiti sulla base dei regolamenti di Ateneo vigenti.
	}
}

\langen{

\cventry{2024 --}{PhD Program Quality Assurance Committee}{\newline{}Member of the Quality Assurance Committee for the PhD Program in ICT}{}{Department of Engineering ``E. Ferrari''}{University of Modena and Reggio Emilia, Italy.}	

\cventry{2020 -- }{Department Quality Assurance Committee}{\newline{}Member of the Quality Assurance Committee}{}{Department of Engineering ``E. Ferrari''}{University of Modena and Reggio Emilia, Italy.}

\cventry{2020 -- 2025}{Master's and Bachelor's Degrees Quality Assurance Committee}{\newline{}Responsible for Quality Assurance for the Master's and Bachelor's Degree Programs in Electronics Engineering}{}{Department of Engineering ``E. Ferrari''}{University of Modena and Reggio Emilia, Italy.}

\cventry{2020 -- 2025}{Internship Committee}{\newline{}Responsible for the Evaluation of Internships}{}{Master's Degree in Electric Vehicle Engineering}{University of Bologna, Italy.}

\cventry{2020 -- 2025}{Didactic Committee}{\newline{}Member of the Committee for the Evaluation and Admission of Candidates}{}{Master's Degree in Electric Vehicle Engineering}{University of Bologna, Italy.}
}

\sectioniten{Partecipazione al collegio dei docenti di dottorati di ricerca accreditati dal Ministero dell'Università e della Ricerca}{Participation in the faculty board of PhD programs recognized by the Ministry}

\langit{
	\cventry{2023 - presente}{Membro di Collegio Docenti di Corsi di Dottorato}{\newline{}Corso di Dottorato di Ricerca in ``Information and Communication Technologies''}{}{Dipartimento di
	Ingegneria ``E. Ferrari''}{Università degli Studi di
	Modena e Reggio Emilia.}

	\cventry{2018 -- 2023}{Membro di Collegio Docenti di Corsi di Dottorato}{\newline{}Corso di Dottorato di Ricerca in Ingegneria
	Meccatronica e dell'Innovazione Meccanica del Prodotto}{}{Dipartimento di
	Tecnica e Gestione dei Sistemi Industriali (DTG)}{Università degli Studi di
	Padova.}

	\cventry{2018 -- 2020}{Coordinamento Didattica Dottorato}{\newline{}Responsabile del Calendario dei Corsi e Seminari dell'area Meccatronica per il Corso di Dottorato di Ricerca in Ingegneria Meccatronica e dell'Innovazione Meccanica del Prodotto}{}{Dipartimento di Tecnica e Gestione dei Sistemi Industriali (DTG)}{Università degli Studi di
	Padova.}
}

\langen{
	\cventry{2023 - presente}{\textbf{Member of the PhD Faculty Board}}{\newline{}PhD Programme in ``Information and Communication Technologies''}{}{Department of Engineering ``E. Ferrari''}{University of Modena and Reggio Emilia, Italy.}

	\cventry{2018 -- 2023}{\textbf{Member of the PhD Faculty Board}}{\newline{}PhD Course in ``Mechatronic and Product Innovation Engineering''}{}{Department of Management and Engineering}{University of Padua, Italy.}

	\cventry{2018 -- 2020}{\textbf{Teaching Coordination within PhD School}}{\newline{}Responsible for the Teaching Coordination within PhD School in ``Mechatronic and Product Innovation Engineering''}{}{Department of Management and Engineering}{University of Padua, Italy.}
}


\subsectioniten{Supervisione di dottorandi nell'ambito di dottorati di ricerca accreditati dal Ministero dell'Università e della Ricerca}{PhD students supervisor}

\langen{
\cvitem{2024 - presente}{PhD Thesis supervisor of Dr. Dario Genzardi, Dept. of Engineering ``E. Ferrari'', University of Modena, Italy}

\cvitem{2020 - 2023}{PhD Thesis supervisor of Dr. Tommaso Fedullo, that is within both the Dept. Management and Engineering, University of Padua, Italy and the Dept. of Engineering ``E. Ferrari'', University of Modena, Italy}

% \cvitem{Co--Supervisor}{--- I have been Co--Supervisor for one former PhD student, Dr. Michele Luvisotto, at the Dept. of Information Engineering, University of Padua, Italy. Dr. Luvisotto successfully defended his PhD thesis and is now within Hitachi PowerGrid Research, Sweden.\medskip\newline
% 	--- I have been  Co--Supervisor of two PhD students, Dr. Alberto Morato and Dr. Giovanni Peserico, both at the Dept. of Information Engineering, University of Padua, Italy, whose Supervisor is prof. S. Vitturi.}

}
\langit{
\cventry{2024 - 2025}{Supervisore di Tesi di Dottorato di Ricerca}{Dr. Dario Genzardi}{Ciclo XXXVIII}{\newline{}Dipartimento di Ingegneria ``E. Ferrari''}{Università di Modena e Reggio Emilia, Italia}

\cventry{2020 - 2023}{Supervisore di Tesi di Dottorato di Ricerca}{Dr. Tommaso Fedullo}{Ciclo XXXV,\newline Dottorato di ricerca svolto presso il Dipartimento di Ingegneria ``E. Ferrari'', Università di Modena, Italia}{\newline{}Dipartimento di Tecnica e Gestione dei Sistemi Industriali}{Università di Padova, Italia}

% \cventry{2019 - 2021}{Co-Supervisore di Tesi di Dottorato di Ricerca}{Dr. Alberto Morato}{Ciclo XXXIV,\newline Dottorato in Alto Apprendistato, svolto presso CMZ Sistemi Elettronici Srl, Treviso, Italia}{\newline{}Dipartimento di Ingegneria dell'Informazione}{Università di Padova, Italia}

% \cventry{2016 - 2018}{Co-Supervisore di Tesi di Dottorato di Ricerca}{Dr. Michele Luvisotto}{Ciclo XXX}{\newline{}Dipartimento di Ingegneria dell'Informazione}{Università di Padova, Italia}%Hitachi PowerGrid Research, Svezia.

%	--- Ho co--supervisionato due dottorandi, Dr. Alberto Morato e Dr. Giovanni Peserico, entrambi presso il Dipartimento di Ingegneria dell'Informazione, Università di Padova, Italia, il cui supervisore è prof. S. Vitturi.}
}



